\section{A finite combinatorial realization}\label{sec:combinatorial}

We construct data satisfying Proposition~\ref{prop:projection-criterion}.
A restricted-intersection bound gives real unit vectors such that every
sufficiently large subfamily contains an orthogonal pair.  The fiber
argument for threefold tensor products strengthens this conclusion to
eight mutually orthogonal members in every sufficiently large subfamily.
Six of these eight vectors will index the six measurements in a parity
obstruction; the remaining two receive a trivial measurement.  Tensoring
these measurements over the eight-subsets gives the projections required
by the criterion.  Because the projections are real, the resulting state
is produced by the same trace-preserving PPT channel on both halves of a
maximally entangled state.

\subsection{Vectors with forced orthogonality}

Set \(p=257\), \(k=4p=1028\), and \([k]=\{1,\ldots,k\}\).  Define
\begin{equation}\label{eq:combinatorial-base}
 \begin{gathered}
 \mathcal U=\{A\subseteq[k]: |A|=2p,\ 1\in A\},\\
 b=|\mathcal U|=\binom{1027}{513}
              =\frac12\binom{4p}{2p},\\
 s_A=\frac{1}{\sqrt k}
       \bigl(2\mathbf1_{\{h\in A\}}-1\bigr)_{h=1}^k
       \in\mathbb R^k
       \quad(A\in\mathcal U).
 \end{gathered}
\end{equation}
Each \(s_A\) is a unit vector.  Since
\(|A\mathbin{\triangle}B|=4p-2|A\cap B|\), counting agreeing and
disagreeing coordinates gives
\begin{equation}\label{eq:combinatorial-inner-product}
 s_A^{\T}s_B
 =\frac{k-2|A\mathbin{\triangle}B|}{k}
 =\frac{|A\cap B|}{p}-1.
\end{equation}
Thus two base vectors are orthogonal exactly when their sets intersect
in \(p\) elements.  Put
\begin{equation}\label{eq:combinatorial-delta}
 a_0=\sum_{j=0}^{p-1}\binom{k}{j}
     =\sum_{j=0}^{256}\binom{1028}{j},
 \qquad \delta=\frac{a_0}{b}.
\end{equation}

Frankl and R\"odl's theory of forbidden intersections also forces
orthogonal families of a prescribed size in large subsets of a sign
cube \cite[Theorem~1.11]{FR1987}.  Here the precise bound follows
from a direct polynomial argument and the fiber amplification below.

\begin{lemma}[Restricted intersections]\label{lem:restricted-intersections}
If \(\mathcal F\subseteq\mathcal U\) contains no two distinct sets
whose associated vectors are orthogonal, then
\(|\mathcal F|\leq a_0=\delta b\).
\end{lemma}

\begin{proof}
The restricted-intersection estimate belongs to the Frankl--Wilson
theory \cite{FW1981}.  The proof below uses the multilinear-polynomial
approach of Alon, Babai, and Suzuki \cite{ABS1991}.  We give the needed
specialization in full.
The integer \(257\) is prime: the primes at most \(\sqrt{257}\) are
\(2,3,5,7,11,13\), and the respective remainders are
\(1,2,2,5,4,10\).

For distinct \(A,B\in\mathcal F\), their common element \(1\) and their
equal sizes give \(1\leq|A\cap B|\leq2p-1\).
Equation~\eqref{eq:combinatorial-inner-product} excludes the value
\(p\).  Hence \(|A\cap B|\) is nonzero modulo \(p\), whereas
\(|A\cap A|=2p\) is zero modulo \(p\).  Over the field \(\F_p\), define
\[
 f_A(x_1,\ldots,x_k)
 =1-\left(\sum_{h\in A}x_h\right)^{p-1}
 \quad(A\in\mathcal F).
\]
For every nonzero \(a\in\F_p\), multiplication by \(a\) permutes the
nonzero field elements; comparing their products gives \(a^{p-1}=1\).
Therefore, on the indicator vectors of the members of \(\mathcal F\),
\[
 f_A(\mathbf1_B)=
 \begin{cases}
  1,&A=B,\\
  0,&A\ne B.
 \end{cases}
\]
Replace every positive power of each variable in every monomial by
its first power.  This preserves all evaluations on Boolean vectors
and leaves degree at most \(p-1\).  The resulting multilinear
polynomials remain linearly independent: evaluating a linear relation
at each \(\mathbf1_B\) makes every coefficient zero.  The space of
multilinear polynomials of degree at most \(p-1\) has a basis consisting
of the \(a_0\) squarefree monomials of those degrees.  Thus
\(|\mathcal F|\leq a_0\).
\end{proof}

\begin{lemma}[Fiber amplification]\label{lem:fiber-amplification}
For \(t=1,2,3\), associate to
\((A_1,\ldots,A_t)\in\mathcal U^t\) the unit vector
\(s_{A_1}\otimes\cdots\otimes s_{A_t}\).
If \(\mathcal R\subseteq\mathcal U^t\) contains no collection of
\(2^t\) distinct tuples whose associated vectors are pairwise
orthogonal, then
\[
 |\mathcal R|\leq t\delta b^t.
\]
\end{lemma}

\begin{proof}
The case \(t=1\) is Lemma~\ref{lem:restricted-intersections}.
For \(t=2\) or \(3\), write
\[
 \mathcal R_A
 =\{(A_2,\ldots,A_t):(A,A_2,\ldots,A_t)\in\mathcal R\},
 \qquad A\in\mathcal U.
\]
Let \(\mathcal G\) consist of those \(A\) for which \(\mathcal R_A\)
contains \(2^{t-1}\) tuples with pairwise orthogonal tensor vectors
in the last \(t-1\) factors.  If distinct \(A,B\in\mathcal G\) had
\(s_A\perp s_B\), choose one such collection in each fiber.
Within either lifted collection, orthogonality comes from the last
factors; between the two collections, it comes from the first factors.
The two fibers are disjoint, so this would produce \(2^t\) distinct
tuples forbidden by the hypothesis.  Thus \(\mathcal G\) has no
orthogonal base pair, and
\(|\mathcal G|\leq a_0=\delta b\) by the preceding lemma.

Each fiber in \(\mathcal G\) has at most \(b^{t-1}\) elements.
Each other fiber has at most \((t-1)\delta b^{t-1}\) elements by the
induction hypothesis.  Summing these bounds and using at most \(b\)
fibers of the latter kind gives
\[
 |\mathcal R|
 \leq (\delta b)b^{t-1}
      +b(t-1)\delta b^{t-1}
 =t\delta b^t.
\]
\end{proof}

Use triples as measurement indices, called \emph{questions}:
\begin{equation}\label{eq:combinatorial-questions}
 \begin{gathered}
 \mathcal Q=\mathcal U^3,\qquad Q=b^3,\qquad
 K=k^3=1028^3=1086373952,\\
 u_q=s_{A_1}\otimes s_{A_2}\otimes s_{A_3}\in\mathbb R^K
 \quad\bigl(q=(A_1,A_2,A_3)\bigr).
 \end{gathered}
\end{equation}
The scalar target is $3\delta<1/K$: the support bound $3\delta Q$
must be strictly below $Q/K$ for the projection criterion.
These vectors are unit vectors.  If \(q=(A_1,A_2,A_3)\) and
\(r=(B_1,B_2,B_3)\), their inner product is the rational number
\begin{equation}\label{eq:combinatorial-question-gram}
 c_{qr}=u_q^{\T}u_r
 =\prod_{h=1}^3\left(\frac{|A_h\cap B_h|}{257}-1\right).
\end{equation}
An \emph{orthogonal eight-set} will mean an eight-element subset of
\(\mathcal Q\) whose associated vectors are pairwise orthogonal.
By Lemma~\ref{lem:fiber-amplification}, with \(t=3\), every
\(\mathcal R\subseteq\mathcal Q\) containing no orthogonal eight-set
satisfies
\begin{equation}\label{eq:combinatorial-eight-bound}
 |\mathcal R|\leq\kappa,\qquad
 \kappa:=3\delta Q=3a_0b^2.
\end{equation}

The next analytic estimate is what makes this bound strict enough for
Proposition~\ref{prop:projection-criterion}.  The central coefficient
is largest among the \(4p+1\) coefficients of \((1+1)^{4p}\), so
\[
 b=\frac12\binom{4p}{2p}\geq\frac{16^p}{2(4p+1)}.
\]
The binomial coefficients increase up to the central one, and the
degree-\(p\) term in the expansion of \((1+1/3)^{4p}\) is no greater
than the full sum.  Hence
\[
 a_0\leq p\binom{4p}{p}
 \leq p(4/3)^{4p}3^p
 =p(256/27)^p
 <p\,12^p.
\]
It follows that
\begin{equation}\label{eq:combinatorial-analytic-bound}
 \delta<2p(4p+1)(3/4)^p<2^{-65}.
\end{equation}
For the last inequality, use \(2p(4p+1)=528906<2^{20}\),
\(p=3\cdot85+2\), and
\((3/4)^3=27/64<1/2\); the remaining factor \((3/4)^2\) is less
than one.  Also \(1028<2^{11}\), so
\[
 K<2^{33},\qquad
 3\delta<3\cdot2^{-65}<2^{-33}<\frac1K.
\]
Combining this with \eqref{eq:combinatorial-eight-bound} proves
\begin{equation}\label{eq:combinatorial-strict-bound}
 \kappa=3a_0b^2<\frac QK.
\end{equation}
The proof has now supplied the vector and cardinality ingredients.
We next construct measurements that turn every orthogonal eight-set
into a forbidden positive-overlap choice.

\subsection{Six parity measurements}

On \(\C^2\), let
\[
 X=\begin{pmatrix}0&1\\1&0\end{pmatrix},\qquad
 Z=\begin{pmatrix}1&0\\0&-1\end{pmatrix},\qquad
 Y=\begin{pmatrix}0&-\mathrm i\\ \mathrm i&0\end{pmatrix}.
\]
Use the Mermin--Peres square
\cite{Mermin1990,Peres1990,Peres1991}; its interpretation in the
setting of quantum magic games is discussed by Arkhipov
\cite{Arkhipov2012}.  The margins in the following array give the
products of the three observables in each row and column:
\begin{equation}\label{eq:combinatorial-parity-square}
 \renewcommand{\arraystretch}{1.25}
 \begin{array}{ccc|c}
  X\otimes I_2&I_2\otimes X&X\otimes X&+I_4\\
  I_2\otimes Z&Z\otimes I_2&Z\otimes Z&+I_4\\
  X\otimes Z&Z\otimes X&Y\otimes Y&+I_4\\\hline
  +I_4&+I_4&-I_4&
 \end{array}
\end{equation}
All nine entries are real symmetric involutions.  In particular, if
\(J=\begin{psmallmatrix}0&-1\\1&0\end{psmallmatrix}\), then
\(Y=\mathrm iJ\) and \(Y\otimes Y=-J\otimes J\), a real integral
symmetric matrix.  Order the six contexts as the three rows followed
by the three columns, each read left to right or top to bottom.  In
context \(t\), write the three observables as \((O_1,O_2,O_3)\).
The first two commute and
\begin{equation}\label{eq:combinatorial-context-signs}
 O_3=e_tO_1O_2,\qquad
 (e_1,e_2,e_3,e_4,e_5,e_6)=(1,1,1,1,1,-1).
\end{equation}
These identities follow from \(XZ=-\mathrm iY\) and
\(ZX=\mathrm iY\).  For example, the first two entries of the last
row multiply, in either order, to \(Y\otimes Y\); those of the last
column multiply, in either order, to \(-Y\otimes Y\).
Since \(O_3\) is a signed product of the first two, all three commute.

Let \(\mathcal A=\{-1,+1\}^2\).  The four outcomes of context \(t\)
are the matrices
\begin{equation}\label{eq:combinatorial-context-projections}
 h_t(\alpha,\beta)
 =\frac14(I_4+\alpha O_1)(I_4+\beta O_2),
 \qquad (\alpha,\beta)\in\mathcal A.
\end{equation}
The two factors divided by two are commuting spectral projections.
Their products are therefore self-adjoint projections.  Distinct
outcomes are orthogonal, and the four products sum to \(I_4\).
They are real symmetric with rational entries, and their ranges have
the three context eigensigns
\((\alpha,\beta,e_t\alpha\beta)\).

\begin{lemma}[Parity obstruction]\label{lem:parity-obstruction}
There is no choice of one projection from each of the six measurements
in \eqref{eq:combinatorial-context-projections} such that the trace
overlap of every two chosen projections from distinct contexts is
strictly positive.
\end{lemma}

\begin{proof}
Suppose there were such a choice.  All six selected projections would
be nonzero.  Let a selected row projection \(E\) and column projection
\(F\) assign eigensigns \(a,b\in\{-1,+1\}\) to their shared
observable \(O\).  The eigenspace identities and self-adjointness give
\[
 aEF=EOF=bEF.
\]
If \(a\ne b\), then \(EF=0\) and their overlap is zero, contrary to
the hypothesis.  Thus the row and column assignments agree at each
of the nine entries of the square.

The choice would consequently give nine consistent signs.  The
product of the three signs in context \(t\) is \(e_t\), by
\eqref{eq:combinatorial-context-signs}.  Multiplying the three row
constraints gives product \(+1\) for the nine signs; multiplying
the three column constraints gives product \(-1\).  This is a
contradiction.  The argument applies to eigenspaces in the complex
Hilbert space \(\C^4\).
\end{proof}

\subsection{Tensor assembly and consequences}

We now impose the six-context obstruction on every orthogonal eight-set
of questions.  Each such set receives its own tensor factor.  Six of its
questions use the six parity measurements on that factor, while
all other questions use the trivial measurement with one outcome $I_4$
and three zero outcomes.  Thus any positive-overlap choice containing
the eight-set violates parity on its factor.  This arrangement also
ensures the commutation needed in the projection criterion: two distinct
questions with nonzero vector inner product cannot both have nontrivial
measurements on any one factor.

For a uniform index set, include a factor carrying only trivial
measurements for every nonorthogonal eight-subset as well.  Let
\begin{equation}\label{eq:combinatorial-tensor-space}
 \begin{gathered}
 \mathcal S=\{S\subseteq\mathcal Q:|S|=8\},\qquad
 L=|\mathcal S|=\binom Q8,\\
 \mathcal H=\bigotimes_{S\in\mathcal S}\C^4\cong\C^D,\qquad
 D=4^L.
 \end{gathered}
\end{equation}
To fix all bases, order subsets of \([k]\) by their increasing lists, then order
tuples lexicographically, and order the eight-subsets by their
increasing lists in that order.  Use the corresponding order for the
tensor factors, the standard tensor-product basis in each
\(\C^4=\C^2\otimes\C^2\), and the
lexicographic order on sign pairs with \(-1<+1\).

For \(S\in\mathcal S\) and \(q\in\mathcal Q\), define a four-outcome
measurement \(h_{S,q}\) on the factor indexed by \(S\).  If \(S\)
is an orthogonal eight-set and \(q\) is its \(t\)-th element with
\(1\leq t\leq6\), set \(h_{S,q}=h_t\).  In every other case, set
\begin{equation}\label{eq:combinatorial-trivial-measurement}
 h_{S,q}(+1,+1)=I_4,\qquad
 h_{S,q}(\alpha,\beta)=0
 \quad\text{when }(\alpha,\beta)\ne(+1,+1).
\end{equation}
Thus the last two questions in an orthogonal eight-set use the
trivial measurement, as do all questions outside that set and all
questions on a nonorthogonal factor.  Every decision of whether a
factor is orthogonal is determined by the rational numbers in
\eqref{eq:combinatorial-question-gram}.

For each \(q\), retain an outcome string on every factor and define
\begin{equation}\label{eq:combinatorial-global-projections}
 \mathcal I_q=\{q\}\times\mathcal A^{\mathcal S},\qquad
 P_{q,\lambda}
   =\bigotimes_{S\in\mathcal S}h_{S,q}(\lambda_S)
   \quad\bigl(\lambda\in\mathcal A^{\mathcal S}\bigr).
\end{equation}
Order these outcome strings lexicographically using the stated factor and
sign-pair orders, and order the full labels first by question and then by
outcome string.  This is the standard basis order used for the concrete maps.
The matrices \(P_{q,\lambda}\) are real symmetric projections.  For a fixed \(q\), two
distinct strings differ in some factor, where their local projections
have zero product; hence their global projections are orthogonal.
Summing over all strings factors as
\(\bigotimes_S\sum_{\alpha,\beta}h_{S,q}(\alpha,\beta)=I_D\).
Keeping every outcome string, including those whose projection is zero,
gives a uniform label set: there are exactly \(4^L=D\) labels for every
question.  Hence
\begin{equation}\label{eq:combinatorial-label-count}
 m:=\left|\bigsqcup_{q\in\mathcal Q}\mathcal I_q\right|=QD.
\end{equation}
In particular \(m\geq D\), as required for the square embedding in
Corollary~\ref{cor:square-embedding}.

\begin{lemma}\label{lem:combinatorial-measurements}
The vectors \eqref{eq:combinatorial-questions} and measurements
\eqref{eq:combinatorial-global-projections} satisfy the two hypotheses
of Proposition~\ref{prop:projection-criterion} with
\(n=m=QD\) and \(\kappa=3a_0b^2<Q/K\).
\end{lemma}

\begin{proof}
If \(q=r\), projections with those labels belong to one projective
measurement and commute.  Suppose that \(q\ne r\) and \(c_{qr}\ne0\).
They cannot both receive nontrivial contexts on any one factor:
that would place them together in an orthogonal eight-set, forcing
\(c_{qr}=0\).  On each factor at least one of the two selected local
projections is consequently \(0\) or \(I_4\).  They commute on each
factor, so their tensor products commute.  This proves the first
hypothesis.

Now let \(J\) be a label set with strictly positive trace overlap
between every distinct pair.  Distinct labels for one question have
zero product, so \(J\) contains at most one label per question.  The
trace overlap of any two global projections factors as
\begin{equation}\label{eq:combinatorial-overlap-product}
 \tr(P_{q,\lambda}P_{r,\mu})
 =
 \prod_{S'\in\mathcal S}
 \tr\bigl(h_{S',q}(\lambda_{S'})
          h_{S',r}(\mu_{S'})\bigr).
\end{equation}
Every local factor is nonnegative by
\eqref{eq:projection-overlap}, whether or not the two local
projections commute.  Thus strict positivity of this finite product
implies strict positivity of each factor.

If the questions represented by \(J\) contained an orthogonal
eight-set \(S\), take the selected labels for its first six questions.
At the factor indexed by \(S\), equation
\eqref{eq:combinatorial-overlap-product} would give strictly positive
overlap for each pair of the six selected context projections.
Lemma~\ref{lem:parity-obstruction} forbids this.  The represented
questions therefore contain no orthogonal eight-set, and
\eqref{eq:combinatorial-eight-bound} gives \(|J|\leq\kappa\).
For completeness, the pairwise condition on a singleton is vacuous
even when its projection is zero; it still satisfies the bound
because \(a_0,b\geq1\) and \(\kappa=3a_0b^2\geq3\).  The strict inequality
\(\kappa<Q/K\) is \eqref{eq:combinatorial-strict-bound}.
\end{proof}

\begin{theorem}[Finite combinatorial construction]\label{thm:old-construction}
Set
\begin{equation}\label{eq:combinatorial-exact-data}
 \begin{gathered}
 p=257,\qquad k=1028,\qquad
 b=\binom{1027}{513},\qquad
 a_0=\sum_{j=0}^{256}\binom{1028}{j},\\
 Q=b^3,\qquad K=1086373952,\qquad
 \kappa=3a_0b^2<\frac QK,\\
 L=\binom Q8,\qquad D=4^L,\qquad m=QD.
 \end{gathered}
\end{equation}
Use the vectors in \eqref{eq:combinatorial-questions}, the \(D\)
labels per question and projections in
\eqref{eq:combinatorial-global-projections}, and the map
\(\Psi:\Mat m\to\Mat D\) in \eqref{eq:projection-map}.
Put \(\Lambda=\Psi^\dagger\) and
\(\Phi_D=\Omega_D/\sqrt D\).
\begin{enumerate}
 \item The map \(\Lambda/Q:\Mat D\to\Mat m\) is a trace-preserving
       PPT channel, and
       \begin{equation}\label{eq:combinatorial-state}
        \rho_{\mathrm{comb}}
        =\left(\frac{\Lambda}{Q}\otimes\frac{\Lambda}{Q}\right)
          (\Phi_D\Phi_D^*)
        =\frac{J(\Psi^\dagger\circ\Psi)}{Q^2D}
       \end{equation}
       is an entangled density operator on
       \(\C^m\otimes\C^m\) belonging to \(\mathscr C\).
       It has \(K_D(\rho_{\mathrm{comb}})=0\) and satisfies the uniform
       unhalved \(1/5\) gap and all operational consequences of
       Theorem~\ref{thm:zero-key} and
       Proposition~\ref{prop:operational-extensions}, with their stated
       protocol, purification, and completion hypotheses.
 \item For the real coordinate isometry \(U:\C^D\to\C^m\) onto the
       first \(D\) coordinates, the maps
       \[
        \Phi_1^{\mathrm{comb}}=\Ad_U\circ\Psi,\qquad
        \Phi_2^{\mathrm{comb}}=\Lambda\circ\Ad_{U^*}
                 =(\Phi_1^{\mathrm{comb}})^\dagger
       \]
       are PPT maps on \(\Mat m\), and
       \begin{equation}\label{eq:combinatorial-pair}
        J(\Phi_2^{\mathrm{comb}}\circ\Phi_1^{\mathrm{comb}})
          =Q^2D\,\rho_{\mathrm{comb}}
       \end{equation}
       is nonseparable over \(\C^m\otimes\C^m\).  In particular their composition is not
       entanglement breaking.
 \item Let \(\Theta_{\mathrm{comb}}\) be the channel obtained by
       applying Proposition~\ref{prop:channel-completion} to this
       square pair.  It is a trace-preserving PPT channel on
       \(\Mat{2m+1}\).  With the positive constants \(c_1,c_2\)
       and the first coordinate inclusion \(W_0\) from that
       proposition, the Choi matrix of its square has the corner
       \begin{equation}\label{eq:combinatorial-channel-corner}
        (W_0^*\otimes W_0^*)J(\Theta_{\mathrm{comb}}^2)
        (W_0\otimes W_0)
        =c_1c_2Q^2D\,\rho_{\mathrm{comb}}.
       \end{equation}
       Consequently
       \(\Theta_{\mathrm{comb}}^2
         =\Theta_{\mathrm{comb}}\circ\Theta_{\mathrm{comb}}\)
       is not entanglement breaking.
\end{enumerate}
\end{theorem}

\begin{proof}
Lemma~\ref{lem:combinatorial-measurements} verifies the hypotheses of
Proposition~\ref{prop:projection-criterion}.  Its composite
\(C=J(\Psi^\dagger\circ\Psi)\) is positive and nonseparable over
\(\C^m\otimes\C^m\), with \(\tr C=Q^2D\).  The projections are real
symmetric, so Corollary~\ref{cor:real-channel} proves
\eqref{eq:combinatorial-state}, the channel assertion, and membership
in \(\mathscr C\).  Theorem~\ref{thm:zero-key} and
Proposition~\ref{prop:operational-extensions} give the stated
operational conclusions.

Since \(m=QD\geq D\), Corollary~\ref{cor:square-embedding} applies to
the displayed coordinate isometry and gives the square PPT pair.
Its composite is \(\Psi^\dagger\circ\Psi\), and the real
identification \eqref{eq:projection-real-identification} gives
\eqref{eq:combinatorial-pair}.  Applying
Proposition~\ref{prop:channel-completion} to this pair gives the
channel on \(\Mat{2m+1}\) and its exact corner
\eqref{eq:combinatorial-channel-corner}.  The corner is nonseparable,
so that proposition also proves that the channel square is not
entanglement breaking.
\end{proof}

All these objects are given by finite exact formulas.  The normalized
sign vectors \(s_A\) have coordinates
\(\pm1/\sqrt{1028}=\pm1/(2\sqrt{257})\), which are irrational because
\(257\) is prime.  The Gram entries of the question vectors in
\eqref{eq:combinatorial-question-gram}, rather than the vector coordinates,
enter the Choi blocks.  Those Gram entries are rational; all nine local
observables are integral, and all local and global projections are
rational.  Consequently every block \(H_{ij}\), every coefficient of
\(\Psi,\Lambda,\Phi_1^{\mathrm{comb}},\Phi_2^{\mathrm{comb}}\), and
every entry of \(\rho_{\mathrm{comb}}\) is rational in the specified
bases.  The completion also has rational coefficients: its positive
scalings and flag matrices in Proposition~\ref{prop:channel-completion}
are formed from the rational square maps using traces and rational
operations.  Appendix~\ref{app:supplementary-checks} describes
supplementary exact checks of the small arithmetic and six-context
calculations.
